\documentclass[a4paper,12pt,ja=standard,lualatex]{bxjsarticle}
\usepackage{luatexja} 
\usepackage{graphicx} % [dvipdfmx] は不要（自動最適化されるため削除）
\usepackage{url}
\usepackage{listings}
\usepackage{makecell}
\usepackage{amsmath, amssymb}
\usepackage{tikz}
\usetikzlibrary{shapes.geometric, arrows.meta, positioning}



\begin{document}
\pagestyle{empty}

\begin{center}
\vspace*{1cm}
\large
{\LARGE 卒業研究報告書}\\
\vspace*{0.8cm}
題目\\
\vspace*{1cm}
{\Huge{音声認識を用いた作業ログ}}\\
{\Huge{自動記録システムの構築と}}\\
{\Huge{業務効率化に関する研究}}\\
\vspace{3mm}

\vspace*{3cm}
指導教員\\
\vspace*{0.3cm}
{\LARGE 森山 真光 教授}\\ 
%\vspace*{-0.3cm}
%\underline{\hspace*{5cm}}\\
\vspace*{3cm}
報告者\\
\vspace*{0.3cm}

{23-1-211-0159}\\
\vspace*{0.3cm}
{\Huge 内山 敦也}\\
% \vspace*{-0.3cm}
% \underline{\hspace*{5cm}}\\
\vspace*{0.5cm}
近畿大学情報学部情報学科\\
\vspace*{2cm}
令和8年6月3日提出\\
\end{center}

\newpage
\normalsize

%\tableofcontents
\clearpage
\begin{center}
{\bf \Large 概要}
\end{center} 
本研究では、音声認識技術を用いて作業内容をリアルタイムに記録し、業務効率化を図るシステムの構築を目的とする。
従来の作業記録は手入力に依存しており、入力の手間や記録漏れといった課題が存在する。
本研究では、マイク入力による音声データを解析し、自動的に作業ログとして記録する仕組みを提案する。
高精度な音声認識（Whisper）と柔軟な意味解析（GPT）を組み合わせることで、従来のシステムでは困難であった「自由発話からの正確な作業ログ分類」を実現し、タスク管理や工数管理など、さまざまな業務への応用および効率的な現場改善への貢献を目指す。

\newpage

\tableofcontents
\newpage
\setcounter{page}{1}
\pagestyle{plain}

\section{序論}\label{sec1}
本研究では、音声認識技術を用いて作業内容をリアルタイムに記録し、業務効率化を図るシステムの構築を目的とする。

\subsection{本研究の背景}
従来の作業記録は手入力に依存しており、入力の手間や記録漏れ、また作業時間の正確な把握が困難であるといった課題が存在する。
特に作業中にキーボード操作などの手入力を挟むことは、本来の業務の手を止めることになり、生産性低下の要因ともなり得る。
近年、深層学習を用いた大規模基盤モデルの登場により、音声認識および自然言語処理の精度が飛躍的に向上している。
これらの実際の研究動向を踏まえ、音声認識を活用したハンズフリーな作業記録環境の実現が求められている。

\subsection{本研究の目的}
本研究の目的は、マイク入力による音声データを解析し、自動的に作業ログとして記録する仕組みを構築することで、業務効率化を図ることである。
高精度な音声認識（Whisper）と柔軟な意味解析（GPT）を組み合わせることで、従来のシステムでは困難であった「自由発話からの正確な作業ログ分類」を実現する。
これにより、ユーザはキーボード操作を行うことなく、音声のみで作業記録を行うことが可能となり、入力コストの大幅な削減と正確な作業時間の把握を目指す。

\subsection{本書の構成}
本書の構成は以下の通りである。

\newpage

\section{関連技術・関連研究}\label{sec2}
本研究に関連する先行研究および技術として、主に「音声認識の業務利用」および「基盤モデル（Whisper、GPT）の発展」の観点から述べる。

\subsection{関連技術}
\subsubsection{Whisper}
音声認識モデルであるWhisper \cite{radford2023whisper} は、大規模かつ多様なデータセットで学習されており、雑音環境下や専門用語に対しても高いロバスト性を持つことが示されている。音声入力から正確にテキストへ変換するための基盤技術として用いる。

\subsubsection{GPT}
大規模言語モデル（LLM）であるGPT-4 \cite{openai2023gpt4} などの登場により、テキストの文脈を深く理解し、非構造化データから意図を抽出することが容易になった。本研究では、テキスト化された発話内容からの意味解析や作業カテゴリの決定に活用する。

\subsection{関連研究}
\subsubsection{音声認識の業務利用と課題}
音声入力を活用した作業記録や業務支援システムは、これまでにも医療現場や保守点検の分野を中心に研究・実用化が進められてきた。
河原 \cite{kawahara2016} は、音声認識の実用化動向において、現場の環境雑音や未知語（システムに登録されていない語彙）への対応が実運用上の課題であると指摘している。
従来のシステムは特定の語彙やテンプレートに依存するものが多く、自由発話から柔軟に作業ログを抽出することは困難であった。

本研究はこれらの実際の研究動向を踏まえ、高精度な音声認識（Whisper）と柔軟な意味解析（GPT）を組み合わせることで、従来のシステムでは困難であった「自由発話からの正確な作業ログ分類」を実現する点に新規性がある。

\newpage

\section{研究内容}\label{sec3}
本研究では、マイクから取得した音声データをリアルタイムに処理し、作業ログとして記録するシステムを構築する。具体的には、音声データをテキストに変換し、その内容から作業開始・終了および作業内容を抽出する。

本システムは以下の構成で動作する。
\begin{itemize}
    \item マイクによる音声入力
    \item 音声認識（Whisper）によるテキスト変換
    \item 自然言語処理（GPT）による意図解析
    \item データベースへの作業ログ保存
\end{itemize}
これにより、ユーザはキーボード操作を行うことなく、音声のみで作業記録を行うことが可能となる。

音声認識を用いた作業ログシステムの実現に向けて、以下の4つの観点から開発・検証を行う。

\subsection{音声認識精度の評価}
音声入力から正確にテキストへ変換できるかを評価する。
\paragraph{目的}
音声認識の精度を確認し、実用可能性を検証する。
\paragraph{処理フロー}
マイク入力 $\rightarrow$ 音声データ取得 $\rightarrow$ テキスト変換（Whisper）

\subsection{作業開始・終了の判定}
テキストから作業の開始・終了を判定する。
\paragraph{目的}
作業時間を自動で計測するための基盤を構築する。
\paragraph{処理フロー}
音声テキスト $\rightarrow$ キーワード判定（開始・終了） $\rightarrow$ 作業状態の更新

\subsection{作業内容の分類}
自然言語処理を用いて作業内容を分類する。
\paragraph{目的}
作業の種類を自動的に整理し、分析可能にする。
\paragraph{処理フロー}
音声テキスト $\rightarrow$ GPTによる意味解析 $\rightarrow$ 作業カテゴリの決定

\subsection{作業ログの蓄積と可視化}
記録されたデータを蓄積し、分析可能な形式で保存する。
\paragraph{目的}
作業時間や傾向を把握し、業務改善に活用する。
\paragraph{処理フロー}
作業データ $\rightarrow$ データベース保存 $\rightarrow$ 可視化・分析

\subsection{利用ツール}
本システムの開発および検証には以下のツール・環境を利用する。
\begin{itemize}
    \item \textbf{Whisper}：音声認識によるテキスト変換
    \item \textbf{Python}：システム実装
    \item \textbf{Flask}：Webアプリケーション化（予定）
\end{itemize}

\newpage

\section{結果・考察}\label{sec4}

\subsection{現時点での開発結果}
現在までに得られているシステムの実装結果は以下の通りである。
\begin{itemize}
    \item 音声入力からテキスト変換までの基本機能は実装済みである。
    \item 簡単なキーワードによる開始・終了判定が可能であることを確認した。
    \item 取得した作業ログの簡易的なデータベース保存機能を確認した。
\end{itemize}

\subsection{考察}
音声入力による作業記録は、従来の手入力に比べて入力にかかるコストや手間を大幅に削減できると考えられる。
また、作業の現場でリアルタイムでの記録が可能となることにより、作業時間の正確な把握が可能になると推察される。
一方で、実際の運用においては、音声認識の誤差や、ユーザの曖昧な発話表現に対する解釈の揺らぎが重要な課題となる。
今後はこれらに対して、GPTのプロンプト調整や、特定の業務コンテキストに合わせたフィルタリング処理などの対応が重要になると考えられる。

\newpage

\section{結論・今後の課題}\label{sec5}

\subsection{結論}
音声認識（Whisper）と自然言語処理（GPT）を組み合わせることで、作業ログ記録の自動化が可能となり、従来の記録に伴う手間や漏れをなくし、業務効率の向上に大きく寄与することが期待される。
また、本システムは汎用的なテキスト・意図解析をベースとしているため、タスク管理や工数管理など、さまざまな形態の業務への応用が可能である。

\subsection{今後の課題}
今後は、構築した基本機能のWebアプリケーション化（Flask）を進めるとともに、以下の課題について検証・実装を行う予定である。
\begin{enumerate}
    \item さまざまな騒音環境下における音声認識精度の評価
    \item 自由発話内における「言い直し」や「曖昧な表現」に対する意図抽出の精度向上
    \item 蓄積されたログデータをダッシュボード等で可視化・分析する仕組みの構築
\end{enumerate}

\newpage

\section*{謝辞}
\addcontentsline{toc}{section}{謝辞}
本研究に取り組むにあたり、ご指導、ご鞭撻を賜りました先生、ならびに実証実験や議論に協力してくださった研究室の皆様に深く感謝いたします。

\newpage

\bibliography{refs}
\bibliographystyle{junsrt}

\newpage
\appendix

\end{document}